\documentclass[11pt,letterpaper]{article}
\usepackage[utf8]{inputenc}
\usepackage{geometry}
\usepackage{setspace}

\geometry{left=1.0in,right=1.0in,top=1.0in,bottom=1.0in}
\doublespacing

\setlength{\parindent}{0pt}
\setlength{\parskip}{12pt plus 2pt minus 2pt}

\begin{document}

\thispagestyle{empty}

\begin{center}
    \large\textbf{IN THE UNITED STATES PATENT AND TRADEMARK OFFICE} \\
\end{center}

\vspace{0.2cm}
\noindent\begin{tabular}{p{2.3in}p{4.0in}}
\textbf{In re Patent Application of:} & Reinaldo Javier Aguiar Marcano \\
\textbf{Application Number:} & 19/560,038 \\
\textbf{Confirmation Number:} & 2359 \\
\textbf{Filing Date:} & March 7, 2026 \\
\textbf{Group Art Unit:} & Pre-Exam / TC 2100 \\
\textbf{Examiner:} & ADAM DEMARAY \\
\textbf{Title of Invention:} & OPTICAL COMPUTING ARCHITECTURE AND SYSTEM \\
\end{tabular}

\vspace{0.6cm}
\begin{center}
    \textbf{\large ABSTRACT OF THE DISCLOSURE}
\end{center}

A novel computer system architecture wherein all major internal and external data buses are implemented using massively parallel optical interconnects. The architecture replaces conventional electrical pathways---including the CPU-to-motherboard socket, the memory bus, and peripheral interconnects---with optical interfaces composed of high-density photon emitter and sensor arrays. This unified optical data fabric eliminates the bottlenecks associated with electrical signaling, such as the CPU I/O limit and the 'memory wall.' The resulting system architecture provides an aggregate data bandwidth in the petabit-per-second (Pbps) range, enabling unprecedented performance for data-intensive applications and fostering new design paradigms such as high-density 3D chip stacking and tightly coupled heterogeneous computing.

\end{document}
